\documentclass[11pt]{article}
\usepackage[margin=1in]{geometry}
\usepackage[T1]{fontenc}
\usepackage[utf8]{inputenc}
\usepackage{lmodern}
\usepackage{amsmath,amssymb,amsthm,mathtools}
\usepackage{booktabs,array,enumitem}
\usepackage{xcolor}
\usepackage[hidelinks]{hyperref}
\usepackage{microtype}

\newcommand{\Logic}{\mathcal{ALCI}_{\mathrm{RCC5}}}
\newcommand{\EQ}{\mathsf{EQ}}
\newcommand{\DR}{\mathsf{DR}}
\newcommand{\PO}{\mathsf{PO}}
\newcommand{\PP}{\mathsf{PP}}
\newcommand{\PPI}{\mathsf{PPI}}
\newcommand{\Base}{\mathsf{B}}
\newcommand{\comp}{\mathsf{comp}}
\newcommand{\Cl}{\mathsf{cl}}
\newcommand{\Mos}{\mathsf{Mos}}
\newcommand{\Q}{\mathcal Q}
\newcommand{\T}{\mathcal T}
\newcommand{\Pos}{\mathsf{Pos}}
\newcommand{\Trip}{\mathsf{Trip}}
\newcommand{\Side}{\mathsf{Side}}
\newcommand{\Reach}{\mathrel{\mathsf{Reach}}}
\newcommand{\eqc}{\equiv_{\EQ}}
\newcommand{\blkeq}{\equiv_{\mathrm{blk}}}
\newcommand{\Ex}{\mathsf{Ex}}
\newcommand{\Mate}{\mathsf{Mate}}
\newcommand{\Stratum}{\mathsf{Stratum}}
\newcommand{\Subsrc}{\mathsf{Subsrc}}
\newcommand{\Vrt}{\mathsf{Vert}}

\title{Formal Review of the Consolidated v2 Generated Split-Forest Proof\
       and the 2026-05-26 Verification Bundle}
\author{Review prepared by GPT-5.5 Pro}
\date{May 26, 2026}

\begin{document}
\sloppy
\maketitle

\begin{abstract}
This memorandum reviews the latest consolidated manuscript
\texttt{repaired\_split\_forest\_no\_automata\_proof\_v2\_consolidated}
and the verification archive \texttt{verification-05-26-2026-latest.zip}.
The new version is a genuine improvement: the verification bundle is now
substantially aligned with the stronger claims, the previous V6 no-mosaic
loophole has been closed in the executable checker, strengthened multiple
upward-$\PP$ tests are present, request-cycle tests are present, and arity-4
mosaic tests are included.  The proof architecture is still credible.
However, the mathematical proof is not yet closed.  The remaining issues are
not merely stylistic; several internal inconsistencies remain between the new
recursive verticalization discipline, the old side-interface lemma, the
support-closure definition, the patchwork proof, and the final enumeration
bound.  I therefore recommend treating the current manuscript as strong
progress, not yet as a finished decidability proof.
\end{abstract}

\tableofcontents

\section{Materials reviewed and reproduction checks}

I reviewed the following uploaded materials:
\begin{itemize}[leftmargin=2em]
  \item \texttt{repaired\_split\_forest\_no\_automata\_proof\_v2\_consolidated.pdf};
  \item \texttt{repaired\_split\_forest\_no\_automata\_proof\_v2\_consolidated.tex};
  \item \texttt{verification-05-26-2026-latest.zip}.
\end{itemize}
The TeX source compiles locally with \texttt{pdflatex} and produces a
36-page PDF.  The only compilation warnings I observed after repeated runs
were harmless hyperref PDF-string warnings.

I also unpacked and reran the verification scripts in
\texttt{verification/python}.  The four main scripts named in the README all
exit with code 0:
\begin{center}
\begin{tabular}{p{0.55\linewidth}cc}
\toprule
Script & Exit code & Result \\
\midrule
\texttt{rcc5\_composition\_check.py} & 0 & PASS \\
\texttt{certificate\_checker.py} & 0 & PASS \\
\texttt{small\_model\_oracle.py} & 0 & PASS \\
\texttt{mosaic\_closure\_search.py} & 0 & PASS \\
\bottomrule
\end{tabular}
\end{center}
The additional script \texttt{wp7\_side\_witness\_stress.py} is present in the
archive, but it does not run from the archive alone: it imports project modules
from \texttt{src/} that are not included in the zip, yielding
\texttt{ModuleNotFoundError: No module named 'alcircc5\_reasoner'}.  This is not
a mathematical refutation, but it means that the manuscript's claim that this
script gives empirical ground truth is not reproducible from the supplied
verification package as distributed.

\section{Executive verdict}

The latest version is a significant improvement over the previous manuscript.
In particular:
\begin{itemize}[leftmargin=2em]
  \item the RCC5 composition table is computationally checked;
  \item the V6 no-mosaic loophole is fixed in the checker;
  \item strengthened two-superpart tests, including an incomparability-forcing
        variant, are present;
  \item request-closure regressions are present;
  \item the proof now explicitly attempts recursive verticalization (M6$'$),
        abstract pair labels $L_{\Q}$, overlap cross-labels, joint equality
        discharge, and uniform cycle representatives.
\end{itemize}
These are the right repairs.  Nevertheless, I would not yet state that the
paper establishes decidability.  Several proof-critical components are still
internally inconsistent or under-proved.  The most important remaining defects
are:
\begin{enumerate}[leftmargin=2em,label=(G\arabic*)]
  \item the old side-interface legality lemma still excludes exactly the
        $\PP/\PPI$-comparable side witnesses that M6$'$ was introduced to admit;
  \item the support-closure definition was weakened to per-position coverage
        (SC4$'$), but the quotient-lifting and patchwork proofs still invoke
        per-triple mosaic coverage from the old SC4;
  \item the overlap/universal-propagation argument does not yet prove that
        cross-mosaic universal restrictions are preserved;
  \item the no-automata regularization section still mixes complete-profile
        repetition and extended-profile repetition, and the uniform-cycle
        substitution lemma needs a stronger compatibility proof;
  \item the final enumeration bound contradicts the earlier side-width
        recurrence and checks only V1--V10, although the certificate validity
        definition now has V1--V14;
  \item the verification package is useful evidence but does not prove the
        support-closed patchwork theorem, and one cited script is not
        self-contained.
\end{enumerate}

\section{Positive progress since the previous round}

\subsection{The verification bundle is now materially stronger}

The earlier mismatch between Claude's memo and the uploaded verification folder
has mostly been corrected.  The latest certificate checker contains 26 test
cases, including:
\begin{itemize}[leftmargin=2em]
  \item the strengthened $C_{AB}^{\mathrm{inc}}$ test;
  \item the single-parent and single-chain negative controls;
  \item pure request-cycle positive and negative controls;
  \item explicit no-mosaic regressions for $\exists\DR.A$ and $\exists\PO.A$;
  \item occurrence-freshness tests separating blocking from semantic equality.
\end{itemize}
All of these pass locally.  This is valuable counterexample-search evidence.

\subsection{Recursive verticalization is the right conceptual repair}

The manuscript correctly identifies the round-4 side-witness problem: two
$\DR/\PO$ side witnesses of the same source can themselves be $\PP/\PPI$-
comparable.  The stress concept
\[
  \exists\DR.(A\sqcap\exists\PP.B)\sqcap\exists\DR.B
\]
can be satisfied with two $\DR$-witnesses $y,z$ of the source and with
$y\PP z$.  A flat sibling-frontier restriction to $\{\DR,\PO\}$ is therefore
too strong.  The proposed recursive verticalization discipline (M6$'$), where
a side witness can open a nested local vertical stratum, is the right direction.

\subsection{The executable V6 defect appears fixed}

In the previous verification review I found that a certificate for
$\exists\DR.A$ with no mosaics was accepted vacuously.  The new checker now
contains explicit negative tests for \texttt{exists DR.A with no mosaics} and
\texttt{exists PO.A with no mosaics}, both rejected.  This is a real improvement.

\section{Remaining proof-critical gaps}

\subsection{The old side-interface lemma contradicts the new M6$'$ discipline}
\label{ssec:side-contradiction}

The most immediate mathematical problem is that the consolidated manuscript
contains both of the following claims.

First, the new M6$'$ discussion says that two side witnesses $y,z$ of the same
source may be $\PP/\PPI$-comparable, and that this is handled by opening a
nested local vertical stratum.  This is the intended round-5 repair.

Second, Lemma \texttt{Side-interface legality} still states that the side
interface $\Side(u)$ is legal with vertical positions only
\[
  \{u\}\cup\{\widehat u:\widehat u\eqc u\},
\]
and all remaining positions are sibling-frontier positions.  It then asserts
that labels between sibling-frontier positions lie in
\[
  \{\DR,\PO,\EQ\},
\]
and explicitly excludes the $\PP/\PPI$ options.

These two statements are incompatible.  The stress concept above is exactly a
case where two side witnesses of the same source may have $y\PP z$.  Under
M6$'$ that pair should be co-vertical inside a nested stratum; under the old
side-interface lemma it is forbidden.  The proof of the lemma even says that a
$\PP$-comparable pair of selected witnesses would be ``folded into the vertical
backbone'', but the definition of $\Side(u)$ does not actually construct the
nested stratum that performs this folding.

This is not only a wording problem.  Later completeness steps still invoke
Lemma \texttt{side-placement} to claim that collected side interfaces satisfy
(M1)--(M6).  If that lemma is the old flat-frontier lemma, it loses satisfiable
models.  If it is meant to be the new recursive lemma, then its definition and
proof must be rewritten.

\paragraph{Required repair.}
Delete or replace the current Lemma \texttt{Side-interface legality}.  The
replacement should define $\Side(u)$ as a finite \emph{stratified} side network,
not as one source/equality cluster plus a flat sibling frontier.  It should prove:
\begin{enumerate}[leftmargin=2em,label=(S\arabic*)]
  \item every $\PP/\PPI$ label among side positions is co-vertical in some
        explicitly named nested stratum;
  \item cross-stratum labels are $\DR/\PO$;
  \item the construction has the stated effective width bound;
  \item the collected stratified side interfaces satisfy the same M6$'$ used
        later in the mosaic definition.
\end{enumerate}

\subsection{SC4$'$ is per-position, but the patchwork proof still uses per-triple coverage}
\label{ssec:sc4}

The manuscript changes support closure to a per-position condition (SC4$'$):
for each abstract position symbol $\pi\in\Pos_Q$, there is a mosaic containing
that position with the right type and stratum.  It then says that per-triple
coverage is not required because V11 enforces the RCC5 composition table on the
abstract triple graph.

However, the next lemmas still rely on the old per-triple version of SC4.  In
particular, Lemma \texttt{Quotient lifting} says that every concrete triple
$(u,v,w)$ is covered by a mosaic $M_{q(u,v,w)}$ ``from (SC4)'', and the proof of
Theorem \texttt{Support-closed mosaic patchwork} likewise says that every triple
of occurrences in a finite prefix is realised in some mosaic by (SC4).

This is an internal contradiction:
\[
  \text{SC4$'$ gives mosaics for individual positions,}
\]
whereas
\[
  \text{the proof uses mosaics for triples.}
\]
V11 may be enough to check RCC5 composition on abstract labels, but then the
patchwork proof must be rewritten to use $L_Q$ directly, not a triple mosaic
that SC4$'$ no longer guarantees.

\paragraph{Required repair.}
Choose one of two routes and make the whole proof consistent.

\begin{description}[leftmargin=2em]
  \item[Route A: restore triple coverage.]  Define support closure so that for
  every abstract triple in $\Trip_Q$ there is a mosaic realising that triple.
  Then Lemma \texttt{Quotient lifting} and Theorem \texttt{patchwork} can keep
  their current form, modulo updated M6$'$ stratification.

  \item[Route B: remove triple-mosaic dependence.]  Keep SC4$'$ per-position
  coverage, but prove a new global-label theorem: every concrete pair label in
  the unfolding is exactly $L_Q(q(u),q(v))$; every concrete triple therefore
  satisfies RCC5 composition by V11; and type propagation/existential support
  are handled by separate V6/V7/V13 clauses.  Under this route, Lemma
  \texttt{Quotient lifting} as currently stated must be deleted.
\end{description}
The current manuscript mixes the two routes.

\subsection{The abstract triple map is still circular}

Definition \texttt{Quotient map from concrete to abstract} says that ordered
concrete triples are sent to elements of $\Trip_Q$ ``by (M1)--(M6), (P1)--(P3),
and the patchwork lemma.''  But the patchwork lemma is later proved using this
same quotient-lifting step.  This makes the argument circular.

A non-circular proof should first define, independently of patchwork, how every
concrete occurrence receives an abstract position symbol and how every concrete
pair label is computed from $L_Q$.  Then V11 can be applied directly to concrete
triples.  Alternatively, if triple mosaics are restored, then $q(u,v,w)$ should
be defined from the concrete unfolding and SC4 should provide a mosaic for it,
without appealing to the theorem being proved.

\subsection{Universal propagation across overlap amalgams is under-proved}

Lemma \texttt{Universal-propagation preservation} is still too compressed.  The
hard case is a cross pair $p\in P_1\setminus P_0$, $q\in P_2\setminus P_0$ with
$L_{12}(p,q)=R$ and $\forall R.D\in\tau(p)$.  The proof says that V7 and a
composition chain through an overlap point force $D\in\tau(q)$.  This is not an
immediate consequence of RCC5 composition.  Composition determines possible
relation labels; it does not by itself propagate type information along a newly
declared cross-pair label.

\paragraph{Required repair.}
The certificate needs an explicit type-safety condition for declared cross-labels:
for every declared overlap cross-pair $(p,q)$ with label $R$, if
$\forall R.D\in\tau(p)$, then $D\in\tau(q)$, and symmetrically using inverse
labels.  Then Lemma \texttt{M4-preserved} becomes immediate.  Without such a
clause, overlap amalgamation can introduce a new $R$-edge that violates a
universal restriction.

\subsection{The request-cycle proof still has profile/extended-profile slippage}

The request-cycle lemma itself is close to the desired no-automata replacement:
it uses half-open cycles and fresh occurrences, and it rejects cycles that carry
unfulfilled vertical existential requests.  This is the right idea.

There is still a mismatch between the lemmas using it.  Lemma
\texttt{Request-fulfilled repetition} uses repeated \emph{extended} profiles
$\widehat\Pi$, including pending-set information.  But Lemma \texttt{Lasso
construction by blocking} says to take the first repetition of the ordinary
complete profile $\Pi(u_i)=\Pi(u_j)$ and then invokes the request-cycle lemma.
That is not sufficient: two occurrences can have the same complete boundary
profile but different pending request sets.  The lasso lemma should repeat
$\widehat\Pi$, not just $\Pi$.

There is also a smaller half-open notation issue.  The proof later says to
replace the tail by fresh copies of $[u_i,u_j]$, while the convention defines
the cycle word as $[u_i,u_j)$.  This should be made consistent throughout.

\paragraph{Required repair.}
State the lasso lemma only for repeated extended profiles:
\[
  \widehat\Pi(u_i)=\widehat\Pi(u_j),
\]
and use the half-open segment $[u_i,u_j)$ everywhere.  Then define the finite
cycle validity check to verify request fulfillment on that half-open word.

\subsection{Uniform global regularization needs a stronger compatibility lemma}

Lemma \texttt{Uniform global regularization} says that all rays hitting the same
extended profile can use one representative cycle word.  This is plausible only
if the extended profile records every boundary component needed for
substitution: context scheme, port pattern, side mosaic interfaces, M6$'$ stratum
forest, equality cluster, request status, and every outward relation summary.
The manuscript asserts this through the substitution lemma, but the proof does
not show that the representative cycle is compatible with every occurrence of
that extended profile in every ray.

This is another place where the proof should be made more mechanical.  The
complete profile should be defined as a finite boundary isomorphism type, and
the substitution lemma should explicitly say that equal complete profiles induce
an isomorphism of all boundary data, including the stratum forest and all named
mosaic positions.  Then the uniform-cycle lemma can rely on that isomorphism.

\subsection{The side-width and enumeration bounds are inconsistent}

The manuscript first says that recursive verticalization requires a side-width
recurrence over modal depth:
\[
  m(0)=1,\qquad
  m(d+1)=1+|\Ex_{\DR/\PO}(u)|(1+k_{\Mate}+k_{\rm bd})
         +|\Ex_{\DR/\PO}(u)|m(d),
\]
and explicitly states that this is not polynomial in general.

Later, in the decision-procedure section, the manuscript reverts to a cubic
bound
\[
  m(n)=1+n^3+3n^2+3n=O(n^3)
\]
for every generated occurrence.  This contradicts the earlier recurrence.  The
abstract and the enumeration bound also speak of a computable double-exponential
bound, while the final proof of Lemma \texttt{enumeration} says that the
certificate validity clauses V1--V10 are checked, omitting the newly added
V11--V14.

\paragraph{Required repair.}
Use one bound consistently.  The safer bound is the recurrence-derived
computable bound, for example
\[
  m_*(C_0)=m(\mathrm{md}(C_0)).
\]
Then define the certificate grammar and enumeration bound using $m_*(C_0)$, not
an unrelated cubic $m(n)$.  Also update every decision-procedure statement from
V1--V10 to V1--V14.

\subsection{The certificate grammar has not caught up with M6$'$}

The mosaic definition now uses a stratum forest
\[
  \mathcal V=(\Stratum,\Subsrc),
\]
but several later pieces still treat the vertical component as a simple subset
$\Vrt\subseteq P$.  Examples include the grammar entry
\[
  \text{Stratum }\Vrt_M \subseteq P
\]
and lemmas that write mosaics as $(P,\tau,L,\Vrt)$ rather than
$(P,\tau,L,\mathcal V)$.  This is not merely notation: a subset $\Vrt$ cannot
encode nested strata, parent-position maps, or the co-vertical relation required
by M6$'$.

\paragraph{Required repair.}
Update the grammar, the finite alphabet count, and every mosaic lemma to use the
full finite stratum-forest object.  The enumeration bound must count finite
stratum forests of size at most the side-width bound.

\section{Assessment of the verification bundle}

The verification evidence is useful and has improved.  It now supports several
finite aspects of the proof architecture:
\begin{itemize}[leftmargin=2em]
  \item the RCC5 composition table is correctly implemented;
  \item the strengthened multiple-superpart examples are included;
  \item occurrence-level equality is distinguished from blocking cycles;
  \item no-mosaic DR/PO witnesses are rejected;
  \item arity-4 atomic mosaic networks and overlap amalgamations are tested.
\end{itemize}
However, the verification bundle remains counterexample-search evidence, not a
formal proof.  In particular:
\begin{enumerate}[leftmargin=2em]
  \item The arity-4 mosaic search does not prove the full support-closed
        patchwork theorem for arbitrary finite prefixes.
  \item The checker does not appear to test the new SC4$'$ versus
        triple-coverage inconsistency identified in Section~\ref{ssec:sc4}.
  \item The side-witness stress script is not self-contained in the supplied
        archive, because it depends on missing \texttt{src/} modules.
  \item The small-model oracle is useful for finite sanity checks, but it cannot
        validate infinite regular witnesses such as
        $\exists\PP.\top\sqcap\forall\PP.\exists\PP.\top$ beyond confirming
        that no small finite model exists.
\end{enumerate}

I recommend including either the full project source needed by
\texttt{wp7\_side\_witness\_stress.py}, or replacing that script with a
self-contained finite model construction for the stress concept.  The latter
would be easy: take three regions/elements $x,y,z$ with $x\DR y$, $x\DR z$, and
$y\PP z$, put $A$ on $y$ and $B$ on $z$, then verify the concept directly.

\section{Recommended next revision}

The next revision should prioritize coherence over adding new machinery.  I
recommend the following concrete edit plan.

\begin{enumerate}[leftmargin=2em,label=(R\arabic*)]
  \item Replace the old \texttt{Side-interface legality} lemma with a genuinely
        recursive M6$'$ side-interface construction.  Do not use the old flat
        sibling-frontier lemma anywhere in completeness.
  \item Decide whether support closure is per-triple or per-position.  If
        per-position, rewrite the patchwork proof to use $L_Q$ directly and
        remove Lemma \texttt{Quotient lifting}.  If per-triple, restore SC4 as
        triple coverage.
  \item Add an explicit cross-pair universal-safety clause for overlap
        cross-labels, and use it to prove Lemma \texttt{M4-preserved}.
  \item Change the lasso construction to repeat extended profiles
        $\widehat\Pi$, and standardize all cycle words as half-open
        $[i,j)$ segments.
  \item Define the complete profile as a finite boundary isomorphism type that
        includes the full stratum forest, all ports, all side mosaic positions,
        and all request statuses.  Then restate the substitution lemma in terms
        of that boundary isomorphism.
  \item Replace the cubic side-width bound by the recurrence-derived computable
        bound $m(\mathrm{md}(C_0))$, or prove that recursive verticalization can
        actually be compressed to the cubic bound.  At present the two claims
        conflict.
  \item Update the certificate grammar and enumeration proof to cover V1--V14
        and to count full stratum forests, not only subsets $\Vrt\subseteq P$.
  \item Make \texttt{wp7\_side\_witness\_stress.py} reproducible from the
        supplied zip, or remove manuscript claims that rely on it.
\end{enumerate}

\section{Final assessment}

The proof has moved in the right direction.  I am more confident now than in the
previous round that the intended split-forest architecture can be made to work:
witness thinning, generated containment covers, occurrence-level equality ports,
recursive side verticalization, support-closed mosaics, and request-fulfilled
blocking cycles are the right ingredients.

But the current consolidated manuscript still does not, as written, establish
decidability.  The largest remaining issue is no longer the original density or
Hasse-diagram concern.  It is the mismatch among three layers that now have to
be aligned precisely:
\[
  \text{recursive side strata (M6$'$)},
  \qquad
  \text{abstract pair/triple support},
  \qquad
  \text{finite regular certificate enumeration}.
\]
Once those layers are made internally consistent, the proof will be much closer
to publication-grade.  Until then, the appropriate verdict is:
\[
  \boxed{\text{Substantial progress, but the decidability proof remains open as a formal manuscript.}}
\]

\end{document}
